

In this section we study Karoubi-localising functors.
We begin in \S\reffour{subsection:cofinality} with a preliminary discussion of the behaviour of $\GW$ and $\L$ under idempotent completions. We show that Schlichting's cofinality theorem for Grothendieck-Witt theory 
holds in our setting 
and deduce the existence of Rothenberg sequences for $\L$-theory. 

Our main result in \S\reffour{subsection:localising-invariants} is Proposition~\reffour{proposition:localising-are-stable}, stating that the $\infty$-category of Karoubi-localising functors valued in an additive $\infty$-category with finite limits is stable. We use this result to construct the Karoubi-Grothendieck-Witt functor $\KGW$ as the unique Karoubi-localising delooping of $\GWspace(-^{\natural})$, and deduce that $\KGW$ is the initial Karoubi-localising functor under $\GW$. We then define the Karoubi $\L$-theory functor $\KL = \KGW^{\bord}$ as the bordification of $\KGW$, and show that it is at the same time also the initial Karoubi-localising functor under $\L$. %
We end this section by a comparison result that describes the fibres of $\GW\to \KGW$ and $\L\to \KL$ in terms of algebraic $\K$-theory.

\subsection{Hermitian cofinality and Rothenberg sequences}
\label{subsection:cofinality}%

We recall from Remark~\reftwo{remark:idem-poincare} that if $(\C,\QF)$ is a Poincar\'e $\infty$-category, then $\QF$ extends (via left Kan extension) uniquely to a Poincar\'e structure on the idempotent completion $\C^\natural$ of $\C$. We therefore, by slight abuse of notation, write $(\C^\natural,\QF)$ for the resulting Poincar\'e $\infty$-category and refer to the association $(\C,\QF) \mapsto (\C^\natural,\QF)$ as idempotent completion of Poincar\'e $\infty$-categories.
We then say that a Poincaré functor $(\C,\QF) \to (\D,\QFD)$ is a Karoubi equivalence if it induces an equivalence on idempotent completions in this sense. Concretely, this means that $\C \to \D$ is a dense inclusion of stable $\infty$-categories, i.e.\ a fully faithful functor such that every object in $\D$ is a retract of an object in $\C$, and $\QFD|_{\C} = \QF$.

Our goal in this subsection is to study the behaviour %
of the Grothendieck-Witt spectrum along Karoubi equivalences $(\C,\QF) \to (\D,\QFD)$.
For this, let us first recall Thomason's cofinality theorem. To ease notation, in the following, for an abelian group $A$ we shall also denote by $A$ its associated Eilenberg-Mac Lane spectrum.

\begin{proposition}
\label{proposition:cofinality-K-theory}%
Let $\C \to \D$ be a Karoubi equivalence of stable $\infty$-categories (e.g.\ $\D = \C^{\natural}$).
Then the induced map 
\[
\K_n(\C)\to \K_n(\D)
\]
is an isomorphism for $n>0$ and a monomorphism for $n=0$. In particular, the square of spectra
\[
\begin{tikzcd}
\K(\C)\ar[r]\ar[d] & \K(\D)\ar[d]\\
\K_0(\C)\ar[r] & \K_0(\D)
\end{tikzcd}
\]
is exact. %
In addition, an object $x\in \D$ is in the essential image of $\C$ %
if and only if its $\K$-class $[x] \in \K_0(\D)$ is in the image of $\K_0(\C)$. 
\end{proposition}
\begin{proof}
The first claim is %
\cite{thomason-trobaugh}*{Theorem 1.10.1} (revisited in the setting of stable $\infty$-categories in \cite{HLS}*{Theorem 7.1}). The last claim follows from Thomason's classification of dense subcategories \cite{thomason-classification}*{Theorem 2.1}. 
\end{proof}



The following result extends~\cite{schlichting-exact}*{Corollary 5.2} to the setting of Poincaré $\infty$-categories.
\begin{proposition}
\label{proposition:weak-cofinality}%
Let $(\C,\QF) \to (\D,\QFD)$ be a Karoubi equivalence of Poincaré $\infty$-categories (e.g.\ $(\D,\QFD) = (\C^\natural,\QF)$).
Then the induced map
\[
\GW_n(\C,\QF)\to \GW_n(\D,\QFD)
\]
is an isomorphism for $n>0$ and a monomorphism for $n=0$, whose image is the subgroup
\[
\GW_0(\D,\QFD) \times_{\K_0(\D)}\K_0(\C)\subseteq \GW_0(\D,\QFD) .
\]
\end{proposition}


\begin{proof}%
We will prove the statement by induction on $n$. Let us first deal with the case $n=0$. By Corollary~\reftwo{corollary:pi-0gw}, the group $\GW_0(\C,\QF)$ is the quotient of the monoid of equivalence classes of Poincaré objects under the relations
\[
[x,q]=[\hyp(y)]
\]
for any Poincar\'e object $(x,q)$ with Lagrangian $y\to x$. In particular, a class $[x,q]$ is 0 if and only if there exist Poincar\'e objects $(x_1,q_1),\,(x_2,q_2)$ with Lagrangians $y_i \to x_i$ for $i=1,2$ such that
\[
(x,q)\oplus (x_1,q_1)\oplus \hyp(y_2)\cong (x_2,q_2)\oplus \hyp(y_1)\,.
\]
To simplify the presentation, let us implicitly identify $\C$ with its essential image in $\D$.
Suppose that $[x,q]$ is in the kernel of the map $\GW_0(\C,\QF)\to \GW_0(\D,\QFD)$. Then we can find $(x_1,q_1)$, $(x_2,q_2)$ and $y_1,y_2$ as above in $\D$. But then there exist $y_1'$ and $y_2'$ in $\D$ such that for $j=1,2$ we have that $y_j'' \coloneq y_j\oplus y_j'$ are in $\C$. Set $(x_j',q_j') = (x_j,q_j) \oplus \hyp(y_j')$. Then $(x_j',q_j')$ is Poincar\'e with Lagrangian $y_j''  \in \C$, which implies that $x_j'$ is in $\C$ as well (since it sits in a fibre sequence between $y_j''$ and its dual). As
\[
(x,q)\oplus (x_1',q_1')\oplus \hyp(y''_2)\cong (x_2',q_2')\oplus \hyp(y_1'')
\]
we have that 
$[x,q]$ vanishes in $\GW_0(\C,\QF)$. We thus conclude that the map $\GW_0(\C,\QF) \to \GW_0(\D,\QFD)$ is injective.
To see that the image is as claimed, suppose $[x,q]\in \GW_0(\D,\QFD)$ is an object such that the $\K_0$-class of $x$ lies in $\K_0(\C)$. Then $x$ is equivalent to an object in $\C$ by Proposition~\reffour{proposition:cofinality-K-theory}. Hence $[x,q]$ is in the image of $\GW_0(\C,\QF) \to \GW_0(\D,\QFD)$.
    
For the induction step, we let $n\geq 1$, and 
contemplate the map of metabolic fibre sequences of Example~\reftwo{example:metabolicfseq}
\[
\begin{tikzcd} 
(\C,\QF\qshift{-1}) \ar[r] \ar[d] & \Met(\C,\QF) \ar[r] \ar[d] & (\C,\QF) \ar[d] \\
(\D,\QFD\qshift{-1}) \ar[r] & \Met(\D,\QFD)\ar[r] & (\D,\QFD) \ .
\end{tikzcd}
\]
By Corollary~\reftwo{corollary:GWHyp},
applying $\GW$ to this diagram produces a diagram of long exact sequences of which a fragment is 
\[
\begin{tikzcd}
[column sep=1em] 
\GW_n(\C,\QF\qshift{-1}) \ar[r]\ar[d] & \K_n(\C)\ar[r]\ar[d,"\cong"] &\GW_n(\C,\QF)\ar[r]\ar[d] & \GW_{n-1}(\C,\QF\qshift{-1}) \ar[r]\ar[d]\arrow[dr, phantom, "\lrcorner", very near start]& \K_{n-1}(\C)\ar[d]\\
\GW_n(\D,\QFD\qshift{-1}) \ar[r] & \K_n(\D)\ar[r] & \GW_n(\D,\QFD)\ar[r] & \GW_{n-1}(\D, \QFD\qshift{-1}) \ar[r]& \K_{n-1}(\D)
\end{tikzcd}
\]
In this diagram, the second vertical arrow is an isomorphism for $n \geq 1$ by Proposition~\reffour{proposition:cofinality-K-theory}. In addition, 
by induction hypothesis, 
the rightmost square is cartesian. A diagram chase then shows that the middle map is surjective. Applying this result to the shifted Poincar\'e structures, we find that also the leftmost map is surjective. Another diagram chase then shows that the middle map is injective.
\end{proof}


We now arrive to the main result of the present subsection, which extends the cofinality theorem of \cite{schlichting-derived}*{Theorem~7.9} 
to the Poincaré categorical setup. In what follows, for an abelian group $A$ with $\Ct$-action, we generally write $A_\hC$, $A^\hC$ and $A^\tC$ for the corresponding construction applied to the Eilenberg-MacLane spectrum of $A$.


\begin{theorem}[Hermitian cofinality]
\label{theorem:strong-cofinality}%
Let $(\C,\QF) \to (\D,\QFD)$ be a Karoubi equivalence of Poincaré $\infty$-categories.
Then all the squares in the following diagrams are cartesian:
\[
\begin{tikzcd}
\GW(\C,\QF)\ar[r]\ar[d] & \GW(\D,\QFD) \ar[d] \\
\K(\C,\QF)^\hC\ar[r]\ar[d] & \K(\D,\QFD)^\hC \ar[d] \\
\K_0(\C,\QF)^\hC\ar[r] & \K_0(\D,\QFD)^\hC
\end{tikzcd}
\qquad
\begin{tikzcd}
\L(\C,\QF)\ar[r]\ar[d] & \L(\D,\QFD)\ar[d] \\
\K(\C,\QF)^\tC\ar[r]\ar[d] & \K(\D,\QFD)^\tC \ar[d] \\
\K_0(\C,\QF)^\tC\ar[r] & \K_0(\D,\QFD)^\tC
\end{tikzcd}
\]
\end{theorem}

\begin{proof}%
The bottom squares are cartesian by Proposition~\reffour{proposition:cofinality-K-theory} and the fact that $(-)^\hC$ and $(-)^\tC$ are exact functors, so it is enough to prove that the top squares are cartesian. Since they are the front and back faces of the cube
\[
\begin{tikzcd}
[column sep=0.5em, row sep=0.7em]
  & \GW(\C,\QF)\ar[rr]\ar[dd] \ar[ld]&& \GW(\D,\QFD)\ar[dd]\ar[ld]\\
\L(\C,\QF)\ar[rr, crossing over]\ar[dd] && \L(\D,\QFD)\\
  & \K(\C,\QF)^\hC\ar[rr]\ar[ld] && \K(\D,\QFD)^\hC\ar[ld]\\
\K(\C,\QF)^\tC\ar[rr] && \K(\D,\QFD)^\tC\ar[uu,<-, crossing over]\\
\end{tikzcd}
\]
whose left and right faces are cartesian, there is a natural equivalence between the total fibres of the two squares under consideration. Let us denote this common total fibre by $T$, viewed as a functor from the $\infty$-category of Karoubi equivalences $(\C,\QF) \to (\D,\QFD)$ to spectra,
and note that the task is to show that $T = 0$. Now since $T$ is the total fibre of the back face, Proposition~\reffour{proposition:weak-cofinality} and Proposition~\reffour{proposition:cofinality-K-theory} imply that $T$ takes values in coconnected spectra. On the other hand, since $T$ is also the total fibre of the front face which consists of bordism invariant functors, we have that $T(-\qshift{n}) = \Sigma^nT$. Together, this shows that $\Sigma^nT$ takes values in coconnective for all integers $n$, and hence that $T=0$ as desired.
\end{proof}

Let us highlight that the claim for L-theory is deduced from the simpler one for Grothendieck-Witt theory via the fundamental exact sequence, contrary to usual flow of information utilised in this paper series.



\begin{remark}
\label{remark:Rothenberg_map}%
On $\pi_n$, the map $\xi \colon \L(\C, \QF)\to \K_0(\C,\QF)^{\tC}$ appearing in Theorem~\reffour{theorem:strong-cofinality}
sends the class of $(x,q)\in \Poinc(\C, \QF\qshift{-n})$ to the class of $x$ in the Tate cohomology group
\[
\widehat \mathrm{H}^{-n}(\Ct; \K_0(\C)) \cong \frac{\{a\in \K_0(\C)\mid \Dual_\QF a = (-1)^n a\}}{\{a+(-1)^n\Dual_\QF a\mid a\in \K_0(\C)\}} ;
\]
indeed, this is true by definition in degree $0$, and in degree $n$ follows by identifying the above map with the map induced on $\pi_0$ by $\L(\C, \QF\qshift{-n})\to \K_0(\C,\QF\qshift{-n})^{\tC}$, where we note that the $\Ct$-action on $\K_0(\C, \QF\qshift{-n})$ differs from that of $\K_0(\C, \QF)$ by multiplication with $(-1)^n$ (e.g., by Theorem~\reftwo{theorem:kr-intro-two}). We denote the resulting $\Ct$-action by $\K_0(\C,\QF)^{(-1)^n}$. Now, the %
periodicity isomorphism
\[
\widehat\mathrm{H}^{-n}(\Ct; \K_0(\C)) = \pi_{n}\K_0(\C,\QF)^{\tC} \cong \pi_0\K_0(\C,\QF\qshift{-n})^{\tC} = \widehat\mathrm{H}^0(\Ct; \K_0(\C)^{(-1)^{n}})
\]
afforded by bordism invariance of $\K_0(-)^{\tC}$ coincides with the standard periodicity isomorphism in Tate cohomology; indeed, by naturality and the classification of dense Poincaré subcategories, it is enough to note that if $\K_0(\C,\QF)$ is generated by an element $[x]$ with dual $(-1)^n[x]$ then $\widehat{\mathrm{H}}^{-n}(\Ct;\K_0(\C,\QF))$ is either $0$ or $\ZZ/2$, and in any case admits no non-identity automorphisms.
\end{remark}

For a Poincaré $\infty$-category $(\C,\QF)$ and a subgroup $c\subseteq \K_0(\C)$ closed under the duality, let us write $\C^c \subseteq \C$ for the dense subcategory spanned by all objects whose K-class lies in $c$, and abbreviate $(\C^c,\QF|_{\C^c})$ to $(\C^c,\QF)$.
We refer to the $\GW$- and $\L$-spectra of $(\C^c,\QF)$ as the $c$-controlled $\GW$- and $\L$-spectra of $(\C,\QF)$. 

\begin{corollary}
\label{cor:rothenbergconc}%
Given a Poincaré $\infty$-category $(\C,\QF)$ and a pair of consecutive subgroups $c \subseteq d \subseteq \K_0(\C)$ closed under the duality, we obtain exact sequences
\[
\GW(\C^c,\QF)\to \GW(\C^d,\QF)\xrightarrow{\xi} (d/c)^\hC \quad\text{and}\quad \L(\C^c,\QF)\to \L(\C^d,\QF)\xrightarrow{\xi} (d/c)^\tC,
\]
\end{corollary}
They translate to long exact sequences relating the corresponding controlled Grothendieck-Witt- and $\L$-groups of $(\C,\QF)$ with group- and Tate-cohomology of $\Ct$ with coefficients in $d/c$, respectively. 
In the case of $\L$-theory, the resulting long exact sequence 
\begin{equation}
\label{equation:rothenberg}%
\dots \to \L_n(\C^c, \QF) \to \L_n(\C^d, \QF) \xrightarrow\xi \widehat{\mathrm{H}}^{-n}(\Ct; d/c) \to \L_{n-1}(\C^c, \QF) \to \dots
\end{equation}
is a generalisation of the \defi{Rothenberg sequence} from~\cite{Ranickiyellowbook}*{Proposition~1.10.1}. Following Remark~\reffour{remark:Rothenberg_map}, the map $\xi$ can be explicitly described as sending $(x,q)\in \Poinc(\C^d, \QF\qshift{-n})$ to the element of $\widehat{\mathrm{H}}^{-n}(\Ct; d/c)$ represented by the $\K$-class of $x$ in $d/c$. 


\begin{example}
\label{example:free-dedekind}%
Let $R$ be a Dedekind ring whose field of fractions $K$ is a global field of characteristic not $2$. Recall from Corollary~\refthree{corollary:compute-L-groups-global-field} that its symmetric $\L$-groups are given by
\[
\L_n(\Dperf(R),\QF^\s_R) = \L_n^\s(R) \cong 
\begin{cases}
\W^\s(R) & \text{ for } n \equiv 0 (4) \\
\displaystyle\mathop{\oplus}\limits_{\fp \in \mathcal{I}} \ZZ/2 & \text{ for } n \equiv 1 (4) \\
0 & \text{ for } n \equiv 2 (4) \\
\Pic(R)/2 & \text{ for } n \equiv 3 (4) \ ,
\end{cases}
\]
where $\mathcal{I}$ is the set of prime ideals of $R$ which contain the ideal $(2)$.
To compare this to $\L(\Dfree(R),\QF^\s_R)$, we first recall that 
\[
\tilde \K_0(R) \coloneq \K_0(\Dperf(R))/\K_0(\Dfree(R))) \cong \ker(\K_0(R) \to \K_0(K)) \cong \Pic(R)
\]
where
the induced $\Ct$-action on $\Pic(R)$ is the sign action since every line bundle is inverse to its dual. The maps into Tate cohomology in the Rothenberg sequence are thus maps
\[
\xi_{2k}\colon \L^{\sym}_{2k}(R) \to \Pic(R)[2] \quad\text{and}\quad \xi_{2k+1}\colon \L^{\sym}_{2k+1}(R) \to \Pic(R)/2
\]
into the $2$-torsion part and the quotient modulo $2$ respectively. 

Let us now unwind what these maps are (under the isomorphism of Corollary~\refthree{corollary:compute-L-groups-global-field}):
\begin{enumerate}
\item
For $r = 3\;(\mathrm{mod}\;4)$ the map $\xi_r$ is the identity $\Pic(R)/2 \to \Pic(R)/2$. To see this, note that the $\L$-group $\L^{\sym}_{-1}(R)$ is generated by the images of the unit $[R/\fp,q^{\uni}]$ of $\L^\sym_{0}(R/\fp)$ under the various inflation functors $\Dperf(R/\fp,\QF^{\sym}_{R/\fp}) \to \Dperf(R,(\QF^{\sym}_R)\qshift{1})$ for all primes $\fp$ (after choosing a uniformiser for each $\fp$). 
These images map
to the class mod $2$ of $[\fp]$ in $\L^{\sym}_{-1}(R) = \Pic(R)/2$, while its underlying reduced $\K$-theory class is $[R/\fp] = -[\fp] \in \K_0(\Dperf(R))/\K_0(\Dfree(R))) = \Pic(R)$. Since $[\fp]=-[\fp] \;\mathrm{mod}\;2$ the claim follows from Remark~\reffour{remark:Rothenberg_map}.
\item
For $r = 1\;(\mathrm{mod}\;4)$ the map $\xi_r\colon \oplus_{\fp \in \I}\ZZ/2 \to \Pic(R)/2$ sends the generator of the $\fp$ summand to the class of $[\fp]$ mod $2$.
Here the $\L$-group $\L_1^{\sym}(R)$ is generated by shifted Poincaré objects in the image of the inflation functor $\Dperf(R/\fp,\QF^{\sym}_{R/\fp}) \to \Dperf(R,(\QF^{\sym}_R)\qshift{1})$ for $R/\fp$ of characteristic $2$; the generator of $\L^\sym_{2}(R/\fp) \cong \ZZ/2$ then maps to the generator of the corresponding $\ZZ/2$ term in $\L^{\sym}(R)_{1} = \oplus_{\fp \in \I}\ZZ/2$, while as above its underlying reduced $\K$-class is $[\fp]=-[\fp] \;\mathrm{mod}\;2$, and we use again Remark~\reffour{remark:Rothenberg_map}.
\item
For $r = 2\;(\mathrm{mod}\;4)$ the source of $\xi_r$ is the 0 group.
\item
For $r = 0\;(\mathrm{mod}\;4)$ we do not have an explicit formula, but we claim that $\xi_r\colon \W^{\sym}(R) \to \Pic(R)[2]$ is surjective. Indeed, any element in $\Pic(R)[2]$ can be represented by an ideal $\mathfrak{a} \subseteq R$ such that $\mathfrak{a}\cdot \mathfrak{a} = (a)$ is principal. Then $\mathfrak{a}$ is in particular a projective $R$-module and we can endow it with the unimodular symmetric bilinear form $b_a\colon \mathfrak{a} \times \mathfrak{a} \to R$ defined by
\[
b_a(x,y) = \frac{xy}{a} .
\]
The Poincaré object $(\mathfrak{a},b_a) \in \Poinc(\Dperf(R),\QF^{\sym}_R)$ then has underlying reduced $\K$-class $[\mathfrak{a}] \in \Pic(R)[2]$, so that the map in question is indeed surjective. 
\end{enumerate}


We hence conclude from the Rothenberg sequence above %
that $\L_n^\s(\Dfree(R)) = 0$ for $n=2,3 \;(\mathrm{mod}\; 4)$. Identifying $\L_{4k}^{\s}(\Dfree(R))$ with the Witt group $\W^{\sym}_{\f}(R)$ of free symmetric $R$-modules (using 
Corollary~\refthree{corollary:d-zero}\;\refthreeitem{item:cor-two}; the statement there is for $\Dperf(R)$ but the same argument works for any dense subcategory thereof)
we find that that the Rothenberg sequence in degrees $n=0,1 \;(\mathrm{mod}\; 4)$ takes the following form:
\[
0 \to \Pic(R)[2] \to \L_{4k+1}^\s(\Dfree(R)) \to \displaystyle\mathop{\oplus}\limits_{\fp \in \mathcal{I}} \ZZ/2 \to \Pic(R)/2 \to\W^{\sym}_{\f}(R) \to \W^{\sym}(R) \to \Pic(R)[2] \to 0 .
\]
In particular, if $2$ is invertible in $R$, i.e.\ if $\I = \emptyset$, then $\L_{4k+1}^\s(\Dfree(R)) \cong \Pic(R)[2]$, and the free Witt group sits in an exact sequence
\[
0 \to \Pic(R)/2 \to \W^{\sym}_{\f}(R) \to \W^{\sym}(R) \to \Pic(R)[2] \to 0.
\]
\end{example}



\subsection{Karoubi-localising invariants}
\label{subsection:localising-invariants}%

Recall from Definition~\reftwo{definition:additive} that for an $\infty$-category $\E$ with finite limits, a functor $\F\colon \Catp\to\E$ is called additive, Verdier-localising, or Karoubi-localising if it is reduced and sends split Poincaré-Verdire squares, Poincaré-Verdier squares, or Poincaré-Karoubi squares, to cartesian squares, respectively.
Also recall from Proposition~\reftwo{proposition:localizing-criterion} that the condition that $\F$ is Karoubi-localising is equivalent to the condition that $\F$ is Verdier-localising and Karoubi invariant, that is, inverts Karoubi equivalences, %
a criterion which we will employ several times in the discussion below.
We will denote by 
\[
\Funkloc(\Catp,\E) \subseteq \Funadd(\Catp,\E)
\]
the full subcategory spanned by the Karoubi-localising functors, and note that this subcategory is closed under limits, and also under colimits if $\E$ is stable.

\begin{example}
The spectrum valued functor $(\C,\QF) \mapsto \fib[\GW(\C,\QF) \to \K(\C,\QF)^{\hC}] = \fib[\L(\C,\QF) \to \K(\C,\QF)^{\tC}]$ is Karoubi-localising: Indeed, it is Karoubi-invariant by Theorem~\reffour{theorem:strong-cofinality} and Verdier-localising as it is the fibre of two Verdier-localising functors.
\end{example}


None of the spectrum valued functors $\K, \GW$ and $\L$ are Karoubi-localising. Indeed, by Proposition~\reffour{proposition:cofinality-K-theory}, Proposition~\reffour{proposition:weak-cofinality} and Theorem~\reffour{theorem:strong-cofinality} none of these functors are Karoubi-invariant. %
For the same reason, their space valued counterparts $\Kspace = \Om^{\infty}\K, \GWspace = \Om^{\infty}\GW$ and $\Lspace = \Om^{\infty}\L$ are also not Karoubi-localising.

One may be tempted to remedy this by instead considering the functors $\K(-^\natural), \GW(-^\natural)$, and $\L(-^\natural)$. Clearly, these functors are now Karoubi-invariant. However, they are not Verdier-localising. To see this, let for instance $R$ be the nodal curve over a field $k$ of characteristic different from 2, i.e.\ the pullback of the evaluation at $0,1$ map $k[t] \to k \times k$ along the diagonal, compare Remark~\reftwo{remark:no-verdier} for the computations about $R$ we are about to use. Then the sequence 
\[
\Dperf(R \textrm{ on } t) \to \Dperf(R[t]) \to \Dperf(R[t^{\pm1}])
\]
is a Karoubi sequence which is not a Verdier sequence, see Example~\reftwo{example:app-Verdier-not-karoubi}, and is hence the idempotent completion of the Verdier sequence obtained by replacing $\Dperf(R[t^{\pm1}])$ with its dense subcategory given by the essential image of the localisation functor. 
Therefore, either of the above three functors is Verdier-localising only if it sends this dense inclusion to an equivalence. By the Bass-Heller-Swan decomposition, this dense inclusion induces the direct summand inclusion $\K_0(R) \to \K_0(R) \oplus \K_{-1}(R)$ on $\K_0$; here, we use that $\N\K_0(R) = 0$ as follows from the Bass-Milnor-Murthy excision sequence for the above mentioned pullback description of the nodal curve upon passing to polynomial rings. The (hermitian) cofinality theorems say that this is the case if and only if $\K_{-1}(R)$ vanishes (after applying $\Ct$-fixed points or the $\Ct$-Tate construction, respectively). But $\K_{-1}(R) \cong \ZZ$, so none of these three terms vanish.

Moreover, from the natural equivalence $\L(\C,\QF\qshift{n}) \simeq \Sig^n \L(\C,\QF)$ and the fact that $(-)\qshift{n}$ preserves both Verdier and Karoubi sequences, we deduce that the space-valued functor $\Lspace(-^\natural)$ is also not Verdier (and hence not Karoubi) localising. In contrast, the space-valued functor $\Kspace(-^\natural)$ \emph{is} Karoubi-localising, see \cite{HLS}*{Corollary 7.2}. Similarly, we find:




\begin{proposition}
\label{example:GWspace_Karoubi_localising}%
The space-valued functor $\GWspace(-^{\natural})$ 
is Karoubi-localising.
\end{proposition}
\begin{proof}
The claim is equivalent to saying that $\GWspace$ itself sends all Poincaré-Karoubi squares
\[
 \begin{tikzcd}
  (\C, \QF) \ar[r] \ar[d] & (\C', \QF') \ar[d]\\
  (\D, \QFD) \ar[r] & (\D', \QFD'),
 \end{tikzcd}
\]
consisting of idempotent complete Poincaré $\infty$-categories to exact squares. 
Given such a square we may factor it into a pasting of two pullback squares
\[
\begin{tikzcd}
  (\C, \QF) \ar[r] \ar[d] & (\C', \QF') \ar[d]\\
  (\D_0, \QFD_0) \ar[r] \ar[d] & (\D'_0, \QFD'_0) \ar[d]\\
  (\D, \QFD) \ar[r] & (\D', \QFD')
 \end{tikzcd}
\]
where $\D_0 \subseteq \D$ and $\D_0'\subseteq \D'$ denote the essential images of $\C$ and $\C'$, respectively, with $\QFD_0 \coloneq \QFD|_{\D_0}$ and $\QFD'_0 \coloneq \QFD'|_{\D'_0}$.
By Corollaries~\reftwo{corollary:criterion-poincare-karoubi-projection} and \reftwo{corollary:criterion-poincare-karoubi-is-verdier} 
the vertical arrows in the bottom square are Karoubi equivalences and the vertical arrows the top square are then Poincaré-Verdier projections.
We apply $\GW$ and then consider the map between vertical cofibres in the lower square: According to Proposition~\reffour{proposition:weak-cofinality}, this is a map of
$0$-truncated spectra, which on $\pi_0$ is the map
\[
\GW_0(\D,\QFD) / [(\GW_0(\D,\QFD) \times_{\K_0(\D)} \K_0(\D_0)] \to \GW_0(\D',\QFD') / [(\GW_0(\D',\QFD') \times_{\K_0(\D')} \K_0(\D_0')]
\]
which is itself injective by Thomason's classification of dense subcategories and the fact that the lower square is cartesian. %
We conclude that the total cofibre of the lower square (still after applying $\GW$) is 
0-truncated, and the same holds for the total cofibre of the outer square, as the upper square is exact upon applying $\GW$. The outer square thus has 
(-2)-truncated total fibre after applying $\GW$, and is hence cartesian after applying $\GWspace = \Om^{\infty}\GW$.
\end{proof}












The following result lifts $\GWspace(-^\natural)$ to a spectrum valued Karoubi-localising functor.    
Its analogue for $\Catx$ was obtained previously %
by Cisinski--Khan \cite{CK}*{Theorem 4.4.3} using a different proof. 
Let us also point out that, ignoring Poincaré structures everywhere, 
the proof of Proposition~\reffour{proposition:localising-are-stable} applies mutatis mutandis to prove its analogue for $\Catx$ (compare \cite{Efimov}*{Proposition 1.64} and \cite{stable-envelope} in which the case of localising invariants of exact $\infty$-categories is treated.)

\begin{proposition}
\label{proposition:localising-are-stable}%
Let $\A$ be an additive $\infty$-category with finite limits. The $\infty$-category $\Funkloc(\Catp,\A)$ of $\A$-valued Karoubi-localising functors is stable. In particular, the functor
\[
\Omega^\infty \circ (-)\colon \Funkloc(\Catp,\Spa(\A))\to \Funkloc(\Catp,\A)
\]
is an equivalence. 
\end{proposition}


The proof of Proposition~\reffour{proposition:localising-are-stable} will make use of the following lemma, which will also be used later in \S\reffour{subsection:multiplicative-GW}.

\begin{lemma}
\label{lemma:tensor-preserves-karoubi-squares}%
For every $(\C,\QF) \in \Catp$, the functor $(\C,\QF)\otimes (-)\colon \Catp \to \Catp$ preserves Poincaré-Karoubi squares.
\end{lemma}
\begin{proof}
By Proposition~\reftwo{lemma:Morita-exact-criterion}\;\reftwoitem{item:preserves-K-equivalences}
we have that $(\C,\QF)\otimes (-)$ preserves Karoubi equivalences, and so it will suffice to show that it sends Poincaré-Verdier squares to Poincaré-Karoubi squares. Now given a Poincaré-Verdier square on the left, consider the associated square on the right:
\[
\begin{tikzcd}
(\D,\QFD) \ar[r]\ar[d] & (\E,\QFE) \ar[d] \\
(\D',\QFD') \ar[r] & (\E',\QFE')
\end{tikzcd}
\quad\quad\quad\quad
\begin{tikzcd}
(\C,\QF)\otimes^{\flat}(\D,\QFD) \ar[r]\ar[d] & (\C,\QF)\otimes^{\flat}(\E,\QFE) \ar[d] \\
(\C,\QF)\otimes(\D',\QFD') \ar[r] & (\C,\QF)\otimes(\E',\QFE'),
\end{tikzcd}
\]
where $(\C,\QF)\otimes^{\flat}(\D,\QFD) := [(\C,\QF)\otimes(\D',\QFD')] \times_{[(\C,\QF)\otimes(\D',\QFD')]^{\natural}} [(\C,\QF)\otimes(\D,\QFD)]^{\natural}$, and similarly for $(\C,\QF)\otimes^{\flat}(\E,\QFE)$. The square on the right is then Karoubi equivalent to the image under $(\C,\QF)\otimes (-)$ of the square on the left, and it will hence suffice to show that the former is a Poincaré-Verdier square.
Now the vertical legs in the square on the left are Poincaré-Verdier projections by assumption; they hence remain so after tensoring with $(\C,\QF)$ by Proposition~\reftwo{lemma:Morita-exact-criterion}\;\reftwoitem{item:preserves-PV-proj}, and become Poincaré-Karoubi projections after further taking idempotent completions. Since Poincaré-Karoubi projections are preserved under base change (Corollary~\reftwo{corollary:PK-projpullbackstable}) we conclude that the vertical legs in the right square are Poincaré-Karoubi projections. But they are actually Poincaré-Verdier projections since they are essentially surjective (their sources contain the tensor with $(\C,\QF)$ of the vertical sources in the left square). At the same time, by Proposition~\reftwo{lemma:Morita-exact-criterion}\;\reftwoitem{item:sends-PV-squares-to-PK-squares} the induced map on vertical fibres in the right square is obtained from the corresponding one in the left square by applying $[(\C,\QF)\otimes (-)]^{\natural}$, and is hence an equivalence. By Proposition~\reftwo{proposition:poincare-verdier-square-char} it will now suffice to show that the right square is adjointable, or equivalently, that the square
\[
\begin{tikzcd}
\Ind\C\otimes\Ind\D \ar[r]\ar[d] & \Ind\C\otimes\Ind\E \ar[d] \\
\Ind\C\otimes \Ind\D' \ar[r] & \Ind\C \otimes \Ind\E'
\end{tikzcd}
\]
obtained by taking inductive completions is both horizontally and vertically right adjontable, where we have used that $\Ind$ inverts Karoubi equivalences and is symmetric monoidal when viewed as a functor from $\Catx$ to $\Pr^\mathrm L$. Now recall that the right adjoints arising after applying $\Ind$ are colimit preserving, and so the resulting adjunctions in $\Pr^\mathrm L$ are internal ones. It will hence suffice to show that tensoring with $\Ind(\C)$ preserves vertically/horizontally internally right adjointable squares in $\Pr^\mathrm L$. Indeed, this is because $\Ind(\C)\otimes (-)$ is a $\Pr^\mathrm L$-linear endomorphism of $\Pr^\mathrm L$, and hence also $\Cat$-linear.
\end{proof}





\begin{proof}[Proof of Proposition~\reffour{proposition:localising-are-stable}]
The full subcategory $\Funkloc(\Catp,\A) \subseteq \Fun(\Catp,\A)$ is closed under limits, and so the loop functor on this category is given by postcomposition with the loop functor on $\A$. The second part of the proposition hence follows from the first. We now proceed to prove that $\Funkloc(\Catp,\A)$ is stable. Since it has finite limits and is pointed, it suffices to show that $\Omega$ is an equivalence, see \cite{HA}*{Proposition~1.4.2.11(3)}. We will produce an explicit inverse to it.

To this end, let $(\Spa^{\omega_1},\hat{\QF}^{\uni})$ be the hermitian $\infty$-category of $\omega_1$-compact spectra equipped with the Poincaré structure $\hat{\QF}^{\uni}$ right Kan extended along $(\Spaf)\op \subseteq (\Spa^{\omega_1})\op$ from the universal Poincaré structure. Then 
\[
\Om^{\infty}\hat{\QF}^{\uni}\colon (\Spa^{\omega_1})\op \to \Sp
\]
preserves countable cofiltered limits,
and so by~\cite{HTT}*{Lemma 6.3.3.6} the $\infty$-category $\catforms(\Spa^{\omega_1},\hat{\QF}^{\uni})$ admits countable filtered colimits, and by a similar reasoning $\Funh((\Spa^{\omega_1},\hat{\QF}^{\uni}),(\Spa^{\omega_1},\hat{\QF}^{\uni}))$ admits countable filtered colimits. Taking the sequential colimit of the functors $\oplus_{\{1,...n\}}\id$ we hence obtain an endofunctor $S\colon (\Spa^{\omega_1},\hat{\QF}^{\uni}) \to (\Spa^{\omega_1},\hat{\QF}^{\uni})$ which satisfies the
Eilenberg swindle (in the semi-additive $\infty$-category $\Cath$), 
that is, $S\oplus \id \simeq S$. Applying the right adjoint $\Pairings\colon \Cath \to \Catp$, we find that the Poincar\'e category $\Latt^{\omega_1}(\Spaf, \QF^{\uni}) \coloneq \Pairings(\Spa^{\omega_1},\hat{\QF}^{\uni})$ supports an Eilenberg swindle in the semi-additive $\infty$-category $\Catp$. For any Karoubi-localising functor $\FF\colon \Catp\to \A$ we therefore have $\FF(\Latt^{\omega_1}(\Spaf, \QF^{\uni}))\simeq 0$ (since $\A$ is assumed additive). Now the composite
\[
(\Spaf,\QF^\uni) \to\Pairings(\Spaf,\QF^\uni)\to \Pairings(\Spa^{\omega_1},\hat{\QF}^{\uni}) = \Latt^{\omega_1}(\Spaf, \QF^{\uni}),
\]
is a Poincaré-Verdier inclusion (see Remark~\refone{remark:unit-fully-faithful}) and the corresponding Poincaré-Verdier quotient %
\[
\Tate^{\omega_1}(\Spaf, (\QF^{\uni})\qshift{1}) \coloneq \Latt^{\omega_1}(\Spaf, \QF^{\uni})/(\Spaf,\QF^\uni)
\]
is an $\omega_1$-bounded variant of the Tate category, see \S\reftwo{subsection:poincare-verdier-sequences}. %
Now for any Poincaré $\infty$-category $(\C,\QF)$ the sequence
\[
(\C,\QF) = (\C,\QF) \otimes (\Spaf,\QF^{\uni}) \to (\C,\QF) \otimes \Latt^{\omega_1}(\Spaf, \QF^{\uni}) \to (\C,\QF) \otimes \Tate^{\omega_1}(\Spaf, (\QF^{\uni})\qshift{1})
\]
is a Poincaré-Karoubi sequence by Proposition~\reftwo{lemma:Morita-exact-criterion}\;\reftwoitem{item:sends-PV-squares-to-PK-squares}.
As $(\C,\QF) \otimes \Latt^{\omega_1}(\Spaf, \QF^{\uni})$ inherits from $\Latt^{\omega_1}(\Spaf, \QF^{\uni})$ an Eilenberg swindle endomorphism
we obtain an equivalence
\[
\FF\simeq \Omega\FF((-) \otimes \Tate^{\omega_1}(\Spaf,(\QF^{\uni})\qshift{1}))
\]
of functors $\Catp\rightarrow \A$. %
By Lemma~\reffour{lemma:tensor-preserves-karoubi-squares} tensoring with $\Tate^{\omega_1}(\Spaf,(\QF^{\uni})\qshift{1})$ preserves Poincar\'e-Karoubi squares, 
and so precomposing with it defines an endofunctor of $\Funkloc(\Catp,\A)$, inverse to $\Omega$.


\end{proof}

















We now arrive to the main definition of this paper, which is afforded by Proposition~\reffour{proposition:localising-are-stable} and Example \reffour{example:GWspace_Karoubi_localising}.


\begin{definition}
\label{definition:KGW}%
We define the \defi{Karoubi-Grothendieck-Witt} spectrum $\KGW\colon \Catp \to \Spa$ to be the essentially unique Karoubi-localising functor whose infinite loop space is $\GWspace(-^\natural)$. 
\end{definition}



Inspecting the proof of Proposition~\reffour{proposition:localising-are-stable} we see that the $\KGW(\C,\QF)$ can be represented by the $\Omega$-spectrum object whose $n$'th space is $\GWspace((\Tate^{\omega_1}(\Spaf,(\QF^{\uni})\qshift{1})^{\otimes n}\otimes (\C,\QF))^\natural)$. This Lattice-Tate construction of $\KGW$ can be viewed as a Poincaré categorical analogue of the cone-suspension ring approach of Karoubi and Schlichting. In fact, the latter is recovered by the present construction in the case of rings, as shown by the following remark. 



\begin{remark}
Recall from~\cite{comparison}*{Theorem B} that for a ring $R$ equipped with an invertible module with involution $M$, the genuine symmetric Grothendieck-Witt space $\GWspace^{\gs}(R,M)$ recovers the classical symmetric Grothendieck-Witt space $\GWspace_{\cl}^{\sym}(R,M)$. 
We claim that this comparison extends to Karoubi-Grothendieck-Witt spectra. To see this, recall the classical construction of the Karoubi-Grothendieck spectrum by Karoubi and Schlichting using cone and suspension rings, see e.g.\ \cite{schlichting-mv}*{\S 10.1}.
Here, the cone ring $C(R)$ of $R$ is the infinite matrix ring $C(R)$ consisting of the $R$-valued $\NN\times \NN$-matrices each of whose rows and each of whose columns has finitely many non-zero entries, and the suspension ring $S(R)$ is its quotient by the idempotent two-sided ideal $I(R) \subseteq C(R)$ consisting of the matrices which globally have finitely many non-zero entries. Then $S(R)$ is an Ore-localisation of $C(R)$ and consequently $S(R)$ and $I(R)$ are flat $C(R)$-modules. In particular, $I(R)$ is also idempotent in the derived sense and $C(R) \to S(R)$ is a derived localisation.
The rings $C(R)$ and $S(R)$ inherit invertible modules with involution $C(M)$ and $S(M)$, respectively, where $C(M)$ is given by $M$-valued $\NN\times\NN$-matrices with the same vanishing properties as $C(R)$, and $S(M) = C(M)/I(M)$ where $I(M)$ is spanned by the $\NN\times \NN$-matrices with the same vanishing properties as $I(R)$, the involution being induced in both cases by transposition followed by the pointwise involution of $M$.
Writing $P = \oplus_{i \in \NN} R$ for the $C(R)$-projective $(C(R),R)$-bimodule of column vectors one then constructs a sequence of duality preserving functors
\[
(\Proj(R),\Dual_M) \xrightarrow{P\otimes_R (-)} (\Proj(C(R)),\Dual_{C(M)}) \longrightarrow (\Proj(S(R)),\Dual_{S(M)}) ,
\]
which one can verify to be a fibre sequence of additive categories with duality, using the fact that $P^{\vee}\otimes_{C(R)} P \cong R$ and $P\otimes_R P^{\vee} \cong I(R)$, where $P^{\vee} := \Hom_{C(R)}(P,C(R))$ can be identified with the $(R,C(R))$-bimodule of row vectors. By~\cite{schlichting-mv}*{Theorem 10.2, Lemma 10.7} this sequence induces a fibre sequence with contractible middle term on classical symmetric $\GW$-spaces, yielding an equivalence 
\[
\GWspace_{\cl}^{\sym}(R,M) \xrightarrow{\sim} \Om\GWspace_{\cl}^{\sym}(S(R),S(M)).
\]
Iterating this construction produces the connecting maps in a spectrum, which we will refer to as the classical symmetric Karoubi-Grothendieck-Witt spectrum $\KGW_{\cl}^{\sym}(R,M)$ of $(R,M)$. To obtain an equivalence with the Karoubi-Grothendieck-Witt spectrum of Definition~\reffour{definition:KGW} we simply add that, in light of Proposition~\refone{proposition:restriction-of-2-excisive-functors}, the construction of the above fibre sequence of additive categories works the same to produce a fibre sequence
\[
\bigl(\Dperf(R),\QF^{\gs}_M\bigr) \xrightarrow{P\otimes_R (-)} \bigl(\Dperf(C(R)),\QF^{\gs}_{C(M)}\bigr) \longrightarrow \bigl(\Dperf(S(R)),\QF^{\gs}_{S(M)}\bigr)
\]
of Poincaré $\infty$-categories. It then follows from Proposition~\reftwo{proposition:left-Kan-extension-specific-structures} that it is in fact a Poincar\'e-Karoubi sequence as one may apply Example~\reftwo{example:involution-induction-compatible}~\reftwoitem{item:ore-compatible} to see that $C(M)$ is compatible with the derived localisation $C(R) \to S(R)$.
Similar to \cite{schlichting-mv}*{Lemma 10.7}, by Eilenberg swindle one can show that $\KGW^\gs(C(R),C(M))$ vanishes. Hence, the above sequence determines an equivalence
\[
\Om^{\infty-1}\KGW^{\gs}(R,M) \simeq \Om^{\infty}\KGW^{\gs}(S(R),S(M)) = \GWspace^{\gs}(S(R),S(M)).
\]
In particular, we obtain an equivalence $\KGW^{\gs}(R,M) \simeq \KGW_{\cl}^{\sym}(R,M)$
as claimed. A similar argument applies for the quadratic and even form parameters.
\end{remark}



Recall the universal properties of $\GWspace = \Om^{\infty}\GW$ and $\GW$ established in \papertwo, namely, $\GWspace$ is the initial group-like additive functor under $\Poinc$ (Observation~\reftwo{observation:univ-GW-space}) and $\GW$ is the initial additive functor to spectra equipped with a map $\Poinc \Rightarrow \Omega^\infty\GW$ (Corollary~\reftwo{corollary:univ-GW-spectrum}). %
Since $\KGW$ is in particular additive the composite transformation $\Poinc \Rightarrow \GWspace\Rightarrow \GWspace(-^\natural) = \Omega^\infty \KGW$ determines a map
\[
\GW\Rightarrow \KGW
\]
which agrees with the transformation $\GWspace\Rightarrow\GWspace(-^\natural)$ on infinite loop spaces.
In particular, by Proposition~\reffour{proposition:weak-cofinality} the map
\[
\pi_n \GW \to \pi_n \KGW
\]
is an isomorphism for $n>0$ and a monomorphim for $n=0$. 

\begin{proposition}
\label{proposition:universal-KGW}%
The transformations $\SS[\Poinc] \Rightarrow \GW\Rightarrow \KGW$ exhibit $\KGW$ as the initial Karoubi-localising functor under $\SS[\Poinc]$ and $\GW$.
\end{proposition}

The proof of Proposition~\reffour{proposition:universal-KGW} will make use of the following lemma:
\begin{lemma}
\label{lemma:idem-invariant}%
Let $\F\colon \Catp \to \E$ be a functor. Then $\F(-^\natural)$ is the initial Karoubi invariant functor under $\F$.
\end{lemma}
\begin{proof}
The inclusion of Karoubi invariant functors inside all functors can be realised as the restriction along the localisation $\Catp \to\Catpi$ of $\Catp$ by the Karoubi equivalences. This localisation is left Bousfield with fully-faithful right adjoint given by the inclusion $\Catpi \subseteq \Catp$ as idempotent complete $\infty$-categories. Pre-composition along this right adjoint hence gives a left adjoint to the inclusion of Karoubi invariant functors, as desired.
\end{proof}


\begin{proof}[Proof of Proposition~\reffour{proposition:universal-KGW}]
The two statements are equivalent since the map $\SS[\Poinc] \Rightarrow \GW$ induces an equivalence on mapping spaces into any additive functor (and in particular into any Karoubi-localising functor) by the universal property of $\GW$. It will hence suffice to show that $\KGW$ is the initial Karoubi-localising functor under $\SS[\Poinc]$. By adjunction, this is equivalent to showing that for every Karoubi-localising functor $\GG \colon \Catp \to \Sp$ the composite functor
\[
\Nat(\KGW,\GG) \to \Nat(\Om^{\infty}\KGW,\Om^{\infty}\GG) \to \Nat(\GWspace,\Om^{\infty}\GG) \to \Nat(\Poinc,\Om^{\infty}\GG)
\]
is an equivalence. Here, we note that for the second and third term, it does not matter if we consider the two functors as additive functors to $\Grp_{\Einf}$ or as group-like additive functors to $\Sps$, since the forgetful functor between these two types of functors is an equivalence. The first map is then an equivalence by Proposition~\reffour{proposition:localising-are-stable} applied with $\A = \Grp_{\Einf}$, the second is an equivalence by Lemma~\reffour{lemma:idem-invariant} since $\Om^{\infty}\KGW = \GWspace(-^{\natural})$, and the third is an equivalence by the universal property of $\GWspace$.
\end{proof}

We write 
\[
\KK \coloneq \KGW^{\hyp} \colon \Catp \to \Fun(\B\Ct, \Sp)
\]
for the functor obtained by pointwise composing $\KGW$ with $\Hyp\colon \Catp\to \Fun(\B\Ct, \Catp)$. By abuse of notation we will also use this symbol for the underlying functor $\Catx \to \Sp$, which agrees with non-connective, or Bass $\K$-theory: By the $\Catx$-analogue of Proposition~\reffour{proposition:localising-are-stable} established by Cisinski--Khan, it suffices to show that their infinite loop spaces agree, and both are equivalent to $\Kspace(-^\natural)$.


\begin{proposition}
\label{proposition:universal-KK}%
The transformations $\SS[\core] \Rightarrow \K \Rightarrow \KK$ of functors $\Catp\to \Spa$ exhibit $\KK$ as the initial Karoubi-localising functor under $\SS[\core]$ and $\K$. The same statement holds with $\Catp$ replaced by $\Catx$.
\end{proposition}
\begin{proof}
As in the proof of Proposition~\reffour{proposition:universal-KGW}, the claim for $\Catp$ follows from Proposition~\reffour{proposition:localising-are-stable}
and Lemma~\reffour{lemma:idem-invariant} using the universal properties of $\K$ and $\Kspace$ (which in this setting follow from the universal properties of $\GW$ and $\GWspace$; indeed, since $\Hyp\colon \Catx \to \Catp$ is both left and right adjoint to the forgetful functor $U\colon \Catp \to \Catx$ by Corollary~\refone{corollary:hyp-is-adjoint} we have that $\Hyp\circ U$ is both left and right adjoint to itself, and since pre-composition with $\Hyp\circ U$ preserves additive functors it also preserves universal additive approximations). The analogous claim %
for $\Catx$ can be obtained by a similar argument using the corresponding statements for $\Catx$ (and in particular \cite{CK}*{Theorem 4.4.3}).
\end{proof}


\medskip 

We now turn to the Karoubi-invariant version of $\L$-theory. Recall from Lemma~\reftwo{lemma:criterion}
that a spectrum-valued additive functor is bordism invariant if and only if it vanishes on hyperbolic (or equivalently metabolic) Poincaré $\infty$-categories, and that the inclusion of spectrum-valued bordism invariant additive functors inside all spectrum-valued additive functors admits both a left and a right adjoint, called bordification and cobordification respectively, see Theorem~\reftwo{theorem:localisation}. The map $\GW \Rightarrow \L$ then exhibits $\L$ as the bordification of $\GW$.

\begin{definition}
We define the \emph{Karoubi L-theory functor} $\KL\colon \Catp\to\Spa$ to be the bordification of $\KGW$.
\end{definition}



\begin{lemma}
\label{lemma:bordification-again-KL}%
Let $\A$ be a (co)complete stable $\infty$-category. Then the (co)bordification %
of a Karoubi-localising functor $\FF\colon \Catp \to \A$ is again Karoubi-localising. In particular, $\KL$ is Karoubi-localising. %
 %
\end{lemma}
\begin{proof}
According to Proposition~\reftwo{proposition:phi-bord} (whose proof carries over verbatim from $\A=\Sp$ to an arbitrary stable $\infty$-category $\A$ equipped with the relevant (co)limits), an explicit model for the bordification of $\FF$ is the cofibre of the canonical map $\FF^\hyp(-)_\hC\to \FF$, and an explicit model for the cobordification of $\FF$ is the fibre of the canonical map $\FF\to \FF^\hyp(-)^\hC$. Since $\Hyp\colon \Catp\to \Catp$ preserves Karoubi sequences, both of them have the required property. %
\end{proof}

By the universal property of $\L$, the composite functor $\SS[\Poinc] \Rightarrow \KGW \Rightarrow \KL$ factors through $\L$.
\begin{proposition}
\label{proposition:universal-KL}%
The square
\[
\begin{tikzcd}
\SS[\Poinc]\ar[r]\ar[d] & \KGW \ar[d] \\
\L \ar[r] & \KL
\end{tikzcd}
\]
exhibits $\KL$ as initial in three manners, namely, as a bordism invariant functor under $\KGW$, as a Karoubi-localising functor under $\L$, and as a Karoubi-localising and bordism invariant functor under $\SS[\Poinc]$.
\end{proposition}

\begin{proof}
The first of these manners is by definition, and the last follows from the first by
the universal property of $\KGW$ from Proposition~\reffour{proposition:universal-KGW}. %
Combining this with the universal property of $\L$ we conclude that the map $\L \Rightarrow \KL$ exhibits $\KL$ as the initial bordism invariant and Karoubi-localising functor under $\L$. Since cobordification preserves Karoubi-localising functors by Lemma~\reffour{lemma:bordification-again-KL} we have that $\KL$ is also initial as a Karoubi-localising functor under $\L$.
\end{proof}




















Applying Corollary~\reftwo{corollary:tate-square} we now find that $\KGW$ and $\KL$ fit into the Karoubi-localising avatar of the fundamental fibre square from Corollary~\reftwo{corollary:tate-square-L}.
\begin{corollary}
\label{corollary:Karoubi-fund-square}%
There is an exact square 
\begin{equation}
\label{equation:fundamental-karoubi}%
 \begin{tikzcd}
  \KGW \ar[r,"\bord"] \ar[d,"\fgt"'] & \KL \ar[d,"\Xi"] \\
  \KK^{\hC} \ar[r,"\mathrm{can}"] & \KK^{\tC},
 \end{tikzcd}
\end{equation}
of Karoubi-localising functors $\Catp \rightarrow \Spa$. 
\end{corollary}


\begin{theorem}[Karoubi cofinality]
\label{theorem:Karoubi-pullback-squares}%
The following two squares are cartesian.
\[
\begin{tikzcd}
\GW(\C,\QF)\ar[r]\ar[d] & \KGW(\C,\QF)\ar[d] & \L(\C) \ar[r] \ar[d] & \KL(\C,\QF) \ar[d] \\
\K(\C,\QF)^\hC\ar[r] & \KK(\C,\QF)^\hC & \K(\C,\QF)^\tC \ar[r] & \KK(\C,\QF)^{\tC}.
\end{tikzcd}
\]
In particular, we have 
\begin{enumerate}
\item $\pi_*\cof(\GW(\C,\QF)[\frac{1}{2}] \to \KGW(\C,\QF)[\frac{1}{2}]) \cong 
\begin{cases}
0 & \text{if } *  >0 \\
\big(\KK_0(\C,\QF) / \K_0(\C,\QF)\big) [\frac{1}{2}]^{\Ct} & \text{if } * =0 \\
\KK_*(\C,\QF)[\frac{1}{2}]^{\Ct} & \text{if } *  <0
\end{cases} $
and
\item the map $\L(\C,\QF)[\frac{1}{2}] \to \KL(\C,\QF)[\frac{1}{2}]$ is an equivalence.
\end{enumerate}
\end{theorem}
\begin{proof}
Arguing as in the proof of Theorem~\reffour{theorem:strong-cofinality}, consider the commutative cube
\[
\begin{tikzcd}
[column sep=0.5em, row sep=0.7em]
  & \GW(\C,\QF)\ar[rr]\ar[dd] \ar[ld]&& \KGW(\C,\QF)\ar[dd]\ar[ld]\\
\L(\C,\QF)\ar[rr, crossing over]\ar[dd] && \KL(\C,\QF)\\
  & \K(\C,\QF)^\hC\ar[rr]\ar[ld] && \KK(\C,\QF)^\hC\ar[ld]\\
\K(\C,\QF)^\tC\ar[rr] && \KK(\C,\QF)^\tC\ar[uu,<-, crossing over]\\
\end{tikzcd}
\]
whose back and front faces are the two squares of the theorem and whose left and right faces are the fundamental fibre squares for $\GW$ and $\KGW$ respectively. %
The total fibres of the back and front faces are hence equivalent, and we may write $T(\C,\QF)$ for the common value. Then $T(-)$ is a bordism invariant functor since the front face consists of such functors and $T(-)$ takes values in $(-1)$-truncated spectra since the square of infinite loop spaces underlying the back face is the square
\[
\begin{tikzcd}
\GWspace(\C,\QF)\ar[r]\ar[d] & \GWspace(\C^{\natural},\QF)\ar[d]  \\
\Kspace(\C,\QF)^\hC\ar[r] & \Kspace(\C^{\natural},\QF)^\hC
\end{tikzcd}
\]
whose horizontal fibres are $(-1)$-truncated by Propositions~\reffour{proposition:weak-cofinality} and~\reffour{proposition:cofinality-K-theory}. We conclude that $T=0$ as desired.

To obtain the ``in particular'', we argue as follows. Put $D=\cof(\K(\C,\QF)\to\KK(\C,\QF))$.
By Proposition~\reffour{proposition:cofinality-K-theory},
$D$ is coconnective, with
$\pi_0 D\cong\KK_0(\C,\QF)/\K_0(\C,\QF)$ and
$\pi_i D\cong\KK_i(\C,\QF)$ for $i<0$.
The two cartesian squares identify the cofibres of
$\GW(\C,\QF)\to\KGW(\C,\QF)$ and
$\L(\C,\QF)\to\KL(\C,\QF)$ with
$D^{\hC}$ and $D^{\tC}$, respectively.
Now we recall that $(-)^{\hC}$ commutes with filtered colimits of uniformly bounded above spectra with $\Ct$-action. For spectra with a single homotopy group, this boils down to the commutation of $\Ct$-group cohomology with filtered colimits, which can for instance be deduced from its identification with shifted $\Ct$-group homology. Inductively over the steps in the Whitehead tower, one obtains the desired result for bounded spectra and by passing to the colimit over the Whitehead tower using \cite{NS}*{Lemma I.2.6}, one obtains the desired result for bounded above spectra.
Since $(-)_{\hC}$ commutes with all colimits, we deduce that the canonical comparison maps 
\[
D^{\hC}[\tfrac{1}{2}]
\longrightarrow (D[\tfrac{1}{2}])^{\hC}
\quad\text{and}\quad
D^{\tC}[\tfrac{1}{2}]
\longrightarrow (D[\tfrac{1}{2}])^{\tC}
\]
are equivalences, so the desired result follows from the homotopy fixed point and Tate spectral sequences, respectively.
\end{proof}




\begin{corollary}
\label{corollary:regular-noetherian}%
Let $A$ be a connective ring spectrum such that $\pi_0(A)$ is regular noetherian. Then the map
$\K(A) \to \KK(A)$ is an equivalence. As a consequence, for any Poincaré structure $\QF$ on $\Modp{A}$ the maps
\[
\GW(\Modp{A},\QF) \to \KGW(\Modp{A},\QF) \quad\text{and}\quad \L(\Modp{A},\QF) \to \KL(\Modp{A},\QF)
\]
are equivalences as well.
\end{corollary}
\begin{proof}
By Theorem \reffour{theorem:Karoubi-pullback-squares} the claim for Grothendieck-Witt- and $\L$-theory follows from the claim for $\K$-theory.
The latter should be considered as well-known, and follows for example from combining the following two claims:
\begin{enumerate}
\item
\label{item:a}%
For every connective ring spectrum $A$ and every $n\geq 0$, the map $\K_{-n}(A) \to \K_{-n}(\pi_0(A))$ is an isomorphism. %
\item
\label{item:b}%
For every regular noetherian ring $R$ and all $n> 0$, one has $\K_{-n}(R) = 0$. %
\end{enumerate}
The first statement above appears for example in \cite{BGT}*{Theorem 9.53}. The second is due to Bass. We include a proof for the convenience of the reader.

To see \reffouritem{item:a}, we note that the Bass-Heller-Swan decomposition (see~\cite{saunier}*{Theorem 1.3} for a categorical version that is applicable in particular to ring spectra, or the next section for a proof in the Poincare setting which equally applies to stable categories) %
implies that the map $\K_{-n}(A) \to \K_{-n}(\pi_0(A))$ is a retract of the map $\K_0(A[\ZZ^n]) \to \K_0(\pi_0(A)[\ZZ^n]) \cong \K_0(\pi_0(A[\ZZ^n]))$, so it suffices to prove the case $n=0$, which follows immediately from $A$ and $\pi_0(A)$ having the same projective modules up to isomorphism, see e.g.\ \cite{Lurie-K-theory}*{Corollary 3}. For \reffouritem{item:b}, again using the Bass-Heller-Swan decomposition
we find that $\K_{-n}(R)$ is a retract of $\K_{-1}(R[\ZZ^{n-1}])$. Since with $R$ being so, also $R[\ZZ^{n-1}]$ is regular noetherian, 
and so it suffices to prove the claim for $n=1$. But another application of Bass-Heller-Swan %
shows that 
\[
\K_{-1}(R) \cong \mathrm{coker}\big[\K_0(R[t]) \oplus \K_0(R[t]) \xrightarrow{t \mapsto (t,t^{-1})} \K_0(R[t^{\pm 1}])\big]
\]
(this is of course Bass' original definition of negative $\K$-groups) and since $R[t]$ is also regular noetherian, 
it suffices to show that $\K_0(R[t]) \to \K_0(R[t^{\pm1}])$ is surjective. This is true more generally for any localisation $S \rightarrow S[f^{-1}]$ of a regular coherent ring $S$ at a central element $f \in S$: it clearly suffices to lift all finitely generated projective modules over $S[f^{-1}]$ to perfect modules over $S$. But for a regular coherent ring all finitely presented module are perfect, so such a lift can be given simply by lifting a presentation and clearing denominators, see \papertwo Example 4.5.2.
\end{proof}



\begin{corollary}
\label{corollary:filtcolimpreskgwspectrum}%
The functors $\KGW$ and $\KL$ preserve filtered colimits.
\end{corollary}
\begin{proof}
This follows from Karoubi cofinality (Theorem~\reffour{theorem:Karoubi-pullback-squares}), using that the functors $[\KK(-)/\K(-)]^{\hC}$ and $[\KK(-)/\K(-)]^{\tC}$ preserve filtered colimits (see the proof of loc.\ cit.) and the earlier established facts that the functors $\GW$ and $\L$ preserve filtered colimits, see Corollaries~\reftwo{corollary:filtcolimpresgwspectrum} and~\reftwo{corollary:Ladd}.
\end{proof}



We close with a simple application of Karoubi cofinality regarding the commutation of Grothendieck-Witt and $\L$-spectra with infinite products. Recall that both $\K, \KK \colon \Catx \rightarrow \Spa$ do commute with arbitrary products; this was originally established by Carlsson in the context of exact categories in \cite{carlsson-infinite} and for stable $\infty$-categories in \cite{kaswing} and with a more conceptual proof again in \cite{Efimov}.

\begin{corollary}
The functors $\GW,\KGW,\L \colon \Catp \rightarrow \Sp$ commute with arbitrary products, and $\KL \colon \Catp \rightarrow \Sp$ commutes with products over families $\C_i$ with $\KK(\C_i)$ uniformly bounded below.
\end{corollary}


\begin{proof}
That $\L_i \colon \Catp \rightarrow \Ab$ commutes with products follows immediately from the cycle description of $\L_i(\C,\QF)$ as Poincar\'e objects in $(\C,\QF\qshift{i})$ modulo metabolics. The claim for $\L$ thus follows from homotopy groups commuting with products. 

From the fundamental fibre sequence $\K(\C)_\hC \rightarrow \GW(\C,\QF) \rightarrow \L(\C,\QF)$ it then follows that $\GW$ commutes with infinite products if $\K(-)_\hC$ does. Dually to the earlier argument on the commutation of $(-)^{\hC}$ with filtered colimits, the homotopy orbits functor $(-)_\hC$ commutes with cofiltered limits, and consequently arbitrary products, of uniformly bounded below collections of spectra. The claim for $\GW$ thus follows from the one for $\K$. 

For $\KGW$ the statement is now a direct consequence of the left cartesian square in Theorem~\reffour{theorem:Karoubi-pullback-squares}. The claim about $\KL$ follows similarly from the right square in Theorem~\reffour{theorem:Karoubi-pullback-squares} since with $(-)_{\hC}$ also $(-)^{\tC}$ commutes with infinite products of uniformly bounded below spectra. 
\end{proof}

\begin{remark}
The functor $\KL \colon \Catp \rightarrow \Sp$ does not commute with arbitrary infinite products as we show next. Using the identification $\KL = \L^{\cotower{-\infty}}$ that we will establish in Proposition~\reffour{proposition:decorated-L} below, a counterexample was first constructed by Winges in \cite{winges-Ltheory}. 

In fact, the Karoubi-localising fundamental fibre square~\eqreffour{equation:fundamental-karoubi}
implies that $\KL(\prod_{i \in I} (\C_i,\QF_i)) \to \prod_{i \in I} \KL(\C_i,\QF_i)$ is an equivalence if and only if the map 
\[
\left[ \prod_{i \in I} \KK(\C_i,\QF_i)\right]^{\tC} \to \prod_{i \in I} \KK(\C_i,\QF_i)^{\tC}
\]
is one. A different (possibly simpler) example than the one provided by Winges where this is not the case is the following: Let $\mathscr{O}_\infty$ denote the infinite Cuntz algebra, a complex $\mathrm C^*$-algebra (and thus in particular a ring with involution) with $\KK(\mathscr{O}_\infty) \simeq \K^\mathrm{top}(\mathscr{O}_\infty) \simeq \KU$ with trivial $\Ct$-action; this equivalence has first been observed by Cortinas--Phillips \cite{CP}, see e.g.\ \cite{BLP}*{\S 5.3} for some general explanations. Considering then a countable product of $(\Dperf(\mathscr{O}_\infty),\QF^\sym_{\mathscr O_\infty})$ we need to show that the canonical map 
\[
\left[ \prod_{\mathbb{N}} \KU\right]^{\tC} \to \prod_{\mathbb{N}} \KU^{\tC}
\]
is not an equivalence. To that end, we recall that $\pi_*(\KU^\tC) \cong \mathbb{Q}_2[t^{\pm 1}]$, see e.g.\ \cite{NS}*{Example I.2.3 (iii)} with $|t|=-2$ and $\beta \cdot t = 2$ where $\beta$ denotes the image of the Bott element in $\pi_2(\KU)$ under the canonical map $\KU \to \KU^{\tC}$. Moreover, we have $\pi_*(\KU^{\hC}) = \KU_*[[t]]/(t^2-2t)$ and for any $\KU$-module $M$, the map $M^{\hC} \to M^{\tC}$ is a localisation at $t$. 
Hence, the target of the above comparison map evaluates to $\prod_{\mathbb{N}} \mathbb{Q}_2[t^{\pm 1}]$ while the source is the localisation of $\prod_{\mathbb{N}} \KU_*[[t]]/(t^2-2t)$ at $t$ and the sequence $\{\frac{1}{2^n} = \frac{1}{\beta^nt^n}\}_{n \in \mathbb{N}} \in \prod_\mathbb{N} \mathbb{Q}_2[t^{\pm 1}]$ does not lie in the image of the comparison map. 
\end{remark}
